\documentclass[11pt,letterpaper]{article}
\usepackage[technical-reference]{cornell-notes}
\usepackage{needspace}

\title{Clause 6: Basics - Cornell Notes}
\author{}
\date{}
\setDocTitle{Clause 6: Basics}
\setDocSubtitle{C++ 2024 Cornell Notes}
\setDocOwner{C++ 2024 Cornell Notes}
\setDocAuthor{}
\setDocDate{}
\setCornellCollection{C++ 2024 Cornell Notes}
\setCornellUnitType{Clause}
\setCornellUnitNumber{6}
\setCornellUnitTitle{Basics}
\setCornellCompanionLabel{C++ Technical Reference}
\hypersetup{pdftitle={Clause 6: Basics - Cornell Notes},pdfsubject={C++ technical study notes},pdfauthor={}}

\begin{document}
\maketitle
\begin{CornellReferenceOrientationBox}
\textbf{Purpose.} Builds the core semantic model for declarations, definitions, scope, lookup, linkage, objects, lifetimes, types, execution, initialization, threading, and program start and termination.
\par\textbf{Role.} Acts as the semantic backbone of the core language.
\par\textbf{Principal dependencies.} Consumes lexical tokens and supplies concepts used by expressions, declarations, classes, templates, and the library.
\par\textbf{Primary audience.} all programmers, compiler implementers, and static-analysis authors.
\par\textbf{Major subject groups.} Preamble; Declarations and definitions; One-definition rule; Scope; Name lookup; Program and linkage; Memory and objects; Types; Program execution.
\end{CornellReferenceOrientationBox}
\section{Learning Objectives}
\begin{itemize}[label=\textcolor{CornellReferenceTeal}{$\blacktriangleright$}]
\item Define the governing ideas of Preamble and state the consequences for a program or implementation.
\item Identify the governing ideas of Declarations and definitions and state the consequences for a program or implementation.
\item Distinguish the governing ideas of One-definition rule and state the consequences for a program or implementation.
\item Explain the governing ideas of Scope and state the consequences for a program or implementation.
\item Trace the governing ideas of Name lookup and state the consequences for a program or implementation.
\item Apply the governing ideas of Program and linkage and state the consequences for a program or implementation.
\item Analyze the governing ideas of Memory and objects and state the consequences for a program or implementation.
\item Evaluate the governing ideas of Types and state the consequences for a program or implementation.
\item Diagnose the governing ideas of Program execution and state the consequences for a program or implementation.
\item Compare the governing ideas of General and state the consequences for a program or implementation.
\item Verify the governing ideas of Point of declaration and state the consequences for a program or implementation.
\item Relate the governing ideas of Block scope and state the consequences for a program or implementation.
\item Classify the governing ideas of Function parameter scope and state the consequences for a program or implementation.
\item Justify the governing ideas of Lambda scope and state the consequences for a program or implementation.
\item Audit the governing ideas of Namespace scope and state the consequences for a program or implementation.
\item Summarize the governing ideas of Class scope and state the consequences for a program or implementation.
\end{itemize}
\section{Normative Structure Map}
\small\begin{longtable}{@{}>{\RaggedRight\arraybackslash}p{0.13\textwidth}>{\RaggedRight\arraybackslash}p{0.27\textwidth}>{\RaggedRight\arraybackslash}p{0.19\textwidth}>{\RaggedRight\arraybackslash}p{0.31\textwidth}@{}}
\toprule\textbf{Subclause} & \textbf{Subject} & \textbf{Rule category} & \textbf{Key dependencies / entities}\\\midrule\endfirsthead
\toprule\textbf{Subclause} & \textbf{Subject} & \textbf{Rule category} & \textbf{Key dependencies / entities}\\\midrule\endhead
\bottomrule\endfoot
6.1 & Preamble & core-language semantics & Consumes lexical tokens and supplies concepts used by expressions, declarations, classes, templates, and the library.\\
6.2 & Declarations and definitions & core-language semantics & Consumes lexical tokens and supplies concepts used by expressions, declarations, classes, templates, and the library.\\
6.3 & One-definition rule & core-language semantics & Consumes lexical tokens and supplies concepts used by expressions, declarations, classes, templates, and the library.\\
6.4 & Scope & core-language semantics & General, Point of declaration, Block scope, Function parameter scope, and related subclauses\\
6.5 & Name lookup & core-language semantics & General, Member name lookup, Unqualified name lookup, Argument-dependent name lookup, and related subclauses\\
6.6 & Program and linkage & core-language semantics & Consumes lexical tokens and supplies concepts used by expressions, declarations, classes, templates, and the library.\\
6.7 & Memory and objects & core-language semantics & Memory model, Object model, Lifetime, Indeterminate values, and related subclauses\\
6.8 & Types & core-language semantics & General, Fundamental types, Optional extended floating-point types, Compound types, and related subclauses\\
6.9 & Program execution & core-language semantics & Sequential execution, Multi-threaded executions and data races, Start and termination\\
\end{longtable}\normalsize
\begin{CornellReferenceNormativeBox}A heading maps the clause; it does not replace the detailed conditions. For each operation, read requirements, effects, result, exceptions, complexity, remarks, and notes with their stated normative force.\end{CornellReferenceNormativeBox}
\subsection*{Rule Classification Guide}
\small\begin{longtable}{@{}>{\RaggedRight\arraybackslash}p{0.25\textwidth}>{\RaggedRight\arraybackslash}p{0.67\textwidth}@{}}\toprule\textbf{Label or category} & \textbf{How to read it}\\\midrule\endfirsthead\toprule\textbf{Label or category} & \textbf{How to read it (continued)}\\\midrule\endhead\bottomrule\endfoot
Well-formed / ill-formed & Formation is governed by syntax and semantic rules; an ill-formed program does not become well-formed because one implementation accepts it.\\
Diagnosable rule & A violation requires at least one diagnostic unless the wording places the case outside the diagnosable set.\\
No diagnostic required & The program can be ill-formed while the implementation has no diagnostic obligation.\\
Undefined behavior & No execution requirements are imposed; translation success does not establish safety or portability.\\
Unspecified behavior & One of the permitted outcomes may be chosen without documentation.\\
Implementation-defined behavior & The implementation chooses and documents the behavior.\\
Conditionally supported & The implementation may omit support but documents unsupported constructs.\\
\end{longtable}\normalsize
\section{Cornell Cue-and-Notes}
\Needspace{8\baselineskip}
\subsection*{6.1 Preamble}
\begin{longtable}{@{}>{\RaggedRight\arraybackslash\columncolor{CornellReferenceCueGray}}p{0.28\textwidth}>{\RaggedRight\arraybackslash}p{0.64\textwidth}@{}}
\rowcolor{CornellReferenceNavy}\color{white}\textbf{Cue / recall prompt} & \color{white}\textbf{Notes}\\\endfirsthead
\rowcolor{CornellReferenceNavy}\color{white}\textbf{Cue / recall prompt} & \color{white}\textbf{Notes (continued)}\\\endhead
\arrayrulecolor{CornellReferenceLineGray}
\textbf{What rule applies to Preamble?}\\[-1mm]{\scriptsize\CornellClauseRef{6.1}} & Frames the terminology, applicability, and relationships needed to interpret Preamble; detailed rules remain in the following subclauses.\\\hline
\end{longtable}
\begin{CornellReferenceSummaryBox}Preamble supplies the core-language semantics portion of this clause. The key review move is to connect its local wording to the clause-wide model, then classify each condition by normative force before relying on it.\end{CornellReferenceSummaryBox}
\Needspace{8\baselineskip}
\subsection*{6.2 Declarations and definitions}
\begin{longtable}{@{}>{\RaggedRight\arraybackslash\columncolor{CornellReferenceCueGray}}p{0.28\textwidth}>{\RaggedRight\arraybackslash}p{0.64\textwidth}@{}}
\rowcolor{CornellReferenceNavy}\color{white}\textbf{Cue / recall prompt} & \color{white}\textbf{Notes}\\\endfirsthead
\rowcolor{CornellReferenceNavy}\color{white}\textbf{Cue / recall prompt} & \color{white}\textbf{Notes (continued)}\\\endhead
\arrayrulecolor{CornellReferenceLineGray}
\textbf{What rule applies to Declarations and definitions?}\\[-1mm]{\scriptsize\CornellClauseRef{6.2}} & Defines the core-language rules for Declarations and definitions, including formation, semantic effect, interactions with type and lookup, and any ill-formed or implementation-dependent cases.\\\hline
\end{longtable}
\begin{CornellReferenceSummaryBox}Declarations and definitions supplies the core-language semantics portion of this clause. The key review move is to connect its local wording to the clause-wide model, then classify each condition by normative force before relying on it.\end{CornellReferenceSummaryBox}
\Needspace{8\baselineskip}
\subsection*{6.3 One-definition rule}
\begin{longtable}{@{}>{\RaggedRight\arraybackslash\columncolor{CornellReferenceCueGray}}p{0.28\textwidth}>{\RaggedRight\arraybackslash}p{0.64\textwidth}@{}}
\rowcolor{CornellReferenceNavy}\color{white}\textbf{Cue / recall prompt} & \color{white}\textbf{Notes}\\\endfirsthead
\rowcolor{CornellReferenceNavy}\color{white}\textbf{Cue / recall prompt} & \color{white}\textbf{Notes (continued)}\\\endhead
\arrayrulecolor{CornellReferenceLineGray}
\textbf{What rule applies to One-definition rule?}\\[-1mm]{\scriptsize\CornellClauseRef{6.3}} & Defines the core-language rules for One-definition rule, including formation, semantic effect, interactions with type and lookup, and any ill-formed or implementation-dependent cases.\\\hline
\end{longtable}
\begin{CornellReferenceSummaryBox}One-definition rule supplies the core-language semantics portion of this clause. The key review move is to connect its local wording to the clause-wide model, then classify each condition by normative force before relying on it.\end{CornellReferenceSummaryBox}
\Needspace{8\baselineskip}
\subsection*{6.4 Scope}
\begin{longtable}{@{}>{\RaggedRight\arraybackslash\columncolor{CornellReferenceCueGray}}p{0.28\textwidth}>{\RaggedRight\arraybackslash}p{0.64\textwidth}@{}}
\rowcolor{CornellReferenceNavy}\color{white}\textbf{Cue / recall prompt} & \color{white}\textbf{Notes}\\\endfirsthead
\rowcolor{CornellReferenceNavy}\color{white}\textbf{Cue / recall prompt} & \color{white}\textbf{Notes (continued)}\\\endhead
\arrayrulecolor{CornellReferenceLineGray}
\textbf{What rule applies to General?}\\[-1mm]{\scriptsize\CornellClauseRef{6.4.1}} & Frames the terminology, applicability, and relationships needed to interpret General; detailed rules remain in the following subclauses.\\\hline
\textbf{What rule applies to Point of declaration?}\\[-1mm]{\scriptsize\CornellClauseRef{6.4.2}} & Defines the core-language rules for Point of declaration, including formation, semantic effect, interactions with type and lookup, and any ill-formed or implementation-dependent cases.\\\hline
\textbf{What rule applies to Block scope?}\\[-1mm]{\scriptsize\CornellClauseRef{6.4.3}} & Defines the blocking and synchronization contract for Block scope; ownership, wake conditions, timeout behavior, and happens-before edges must be analyzed together.\\\hline
\textbf{What rule applies to Function parameter scope?}\\[-1mm]{\scriptsize\CornellClauseRef{6.4.4}} & Locates names and entity identity for Function parameter scope; distinguish where a name is visible from whether declarations denote the same entity across contexts.\\\hline
\textbf{What rule applies to Lambda scope?}\\[-1mm]{\scriptsize\CornellClauseRef{6.4.5}} & Locates names and entity identity for Lambda scope; distinguish where a name is visible from whether declarations denote the same entity across contexts.\\\hline
\textbf{What rule applies to Namespace scope?}\\[-1mm]{\scriptsize\CornellClauseRef{6.4.6}} & Locates names and entity identity for Namespace scope; distinguish where a name is visible from whether declarations denote the same entity across contexts.\\\hline
\textbf{What rule applies to Class scope?}\\[-1mm]{\scriptsize\CornellClauseRef{6.4.7}} & Locates names and entity identity for Class scope; distinguish where a name is visible from whether declarations denote the same entity across contexts.\\\hline
\textbf{What rule applies to Enumeration scope?}\\[-1mm]{\scriptsize\CornellClauseRef{6.4.8}} & Locates names and entity identity for Enumeration scope; distinguish where a name is visible from whether declarations denote the same entity across contexts.\\\hline
\textbf{What rule applies to Template parameter scope?}\\[-1mm]{\scriptsize\CornellClauseRef{6.4.9}} & Locates names and entity identity for Template parameter scope; distinguish where a name is visible from whether declarations denote the same entity across contexts.\\\hline
\end{longtable}
\begin{CornellReferenceSummaryBox}Scope supplies the core-language semantics portion of this clause. The key review move is to connect its local wording to the clause-wide model, then classify each condition by normative force before relying on it.\end{CornellReferenceSummaryBox}
\Needspace{8\baselineskip}
\subsection*{6.5 Name lookup}
\begin{longtable}{@{}>{\RaggedRight\arraybackslash\columncolor{CornellReferenceCueGray}}p{0.28\textwidth}>{\RaggedRight\arraybackslash}p{0.64\textwidth}@{}}
\rowcolor{CornellReferenceNavy}\color{white}\textbf{Cue / recall prompt} & \color{white}\textbf{Notes}\\\endfirsthead
\rowcolor{CornellReferenceNavy}\color{white}\textbf{Cue / recall prompt} & \color{white}\textbf{Notes (continued)}\\\endhead
\arrayrulecolor{CornellReferenceLineGray}
\textbf{What rule applies to General?}\\[-1mm]{\scriptsize\CornellClauseRef{6.5.1}} & Frames the terminology, applicability, and relationships needed to interpret General; detailed rules remain in the following subclauses.\\\hline
\textbf{Which declarations does Member name lookup find?}\\[-1mm]{\scriptsize\CornellClauseRef{6.5.2}} & Defines which declarations Member name lookup can find, the scopes searched, and the point at which lookup occurs; access and overload selection are separate later questions.\\\hline
\textbf{Which declarations does Unqualified name lookup find?}\\[-1mm]{\scriptsize\CornellClauseRef{6.5.3}} & Defines which declarations Unqualified name lookup can find, the scopes searched, and the point at which lookup occurs; access and overload selection are separate later questions.\\\hline
\textbf{Which declarations does Argument-dependent name lookup find?}\\[-1mm]{\scriptsize\CornellClauseRef{6.5.4}} & Defines which declarations Argument-dependent name lookup can find, the scopes searched, and the point at which lookup occurs; access and overload selection are separate later questions.\\\hline
\textbf{Which declarations does Qualified name lookup find?}\\[-1mm]{\scriptsize\CornellClauseRef{6.5.5}} & Defines which declarations Qualified name lookup can find, the scopes searched, and the point at which lookup occurs; access and overload selection are separate later questions.\\\hline
\textbf{What rule applies to Elaborated type specifiers?}\\[-1mm]{\scriptsize\CornellClauseRef{6.5.6}} & Defines the core-language rules for Elaborated type specifiers, including formation, semantic effect, interactions with type and lookup, and any ill-formed or implementation-dependent cases.\\\hline
\textbf{What rule applies to Using-directives and namespace aliases?}\\[-1mm]{\scriptsize\CornellClauseRef{6.5.7}} & Defines the core-language rules for Using-directives and namespace aliases, including formation, semantic effect, interactions with type and lookup, and any ill-formed or implementation-dependent cases.\\\hline
\end{longtable}
\begin{CornellReferenceSummaryBox}Name lookup supplies the core-language semantics portion of this clause. The key review move is to connect its local wording to the clause-wide model, then classify each condition by normative force before relying on it.\end{CornellReferenceSummaryBox}
\Needspace{5\baselineskip}\paragraph{Detailed coverage ledger.}
\begin{description}[style=nextline,leftmargin=2.8cm,font=\normalfont\bfseries\color{CornellReferenceTeal}]
\item[6.5.5.1] \textbf{General.} Frames the terminology, applicability, and relationships needed to interpret General; detailed rules remain in the following subclauses.
\item[6.5.5.2] \textbf{Class members.} Defines the core-language rules for Class members, including formation, semantic effect, interactions with type and lookup, and any ill-formed or implementation-dependent cases.
\item[6.5.5.3] \textbf{Namespace members.} Defines the core-language rules for Namespace members, including formation, semantic effect, interactions with type and lookup, and any ill-formed or implementation-dependent cases.
\end{description}
\Needspace{8\baselineskip}
\subsection*{6.6 Program and linkage}
\begin{longtable}{@{}>{\RaggedRight\arraybackslash\columncolor{CornellReferenceCueGray}}p{0.28\textwidth}>{\RaggedRight\arraybackslash}p{0.64\textwidth}@{}}
\rowcolor{CornellReferenceNavy}\color{white}\textbf{Cue / recall prompt} & \color{white}\textbf{Notes}\\\endfirsthead
\rowcolor{CornellReferenceNavy}\color{white}\textbf{Cue / recall prompt} & \color{white}\textbf{Notes (continued)}\\\endhead
\arrayrulecolor{CornellReferenceLineGray}
\textbf{What rule applies to Program and linkage?}\\[-1mm]{\scriptsize\CornellClauseRef{6.6}} & Locates names and entity identity for Program and linkage; distinguish where a name is visible from whether declarations denote the same entity across contexts.\\\hline
\end{longtable}
\begin{CornellReferenceSummaryBox}Program and linkage supplies the core-language semantics portion of this clause. The key review move is to connect its local wording to the clause-wide model, then classify each condition by normative force before relying on it.\end{CornellReferenceSummaryBox}
\Needspace{8\baselineskip}
\subsection*{6.7 Memory and objects}
\begin{longtable}{@{}>{\RaggedRight\arraybackslash\columncolor{CornellReferenceCueGray}}p{0.28\textwidth}>{\RaggedRight\arraybackslash}p{0.64\textwidth}@{}}
\rowcolor{CornellReferenceNavy}\color{white}\textbf{Cue / recall prompt} & \color{white}\textbf{Notes}\\\endfirsthead
\rowcolor{CornellReferenceNavy}\color{white}\textbf{Cue / recall prompt} & \color{white}\textbf{Notes (continued)}\\\endhead
\arrayrulecolor{CornellReferenceLineGray}
\textbf{What rule applies to Memory model?}\\[-1mm]{\scriptsize\CornellClauseRef{6.7.1}} & Defines the core-language rules for Memory model, including formation, semantic effect, interactions with type and lookup, and any ill-formed or implementation-dependent cases.\\\hline
\textbf{What rule applies to Object model?}\\[-1mm]{\scriptsize\CornellClauseRef{6.7.2}} & Defines the core-language rules for Object model, including formation, semantic effect, interactions with type and lookup, and any ill-formed or implementation-dependent cases.\\\hline
\textbf{When does Lifetime begin or end?}\\[-1mm]{\scriptsize\CornellClauseRef{6.7.3}} & Tracks storage and object usability for Lifetime; storage availability alone does not prove that an object's lifetime has begun or remains active.\\\hline
\textbf{What rule applies to Indeterminate values?}\\[-1mm]{\scriptsize\CornellClauseRef{6.7.4}} & Defines the core-language rules for Indeterminate values, including formation, semantic effect, interactions with type and lookup, and any ill-formed or implementation-dependent cases.\\\hline
\textbf{What rule applies to Storage duration?}\\[-1mm]{\scriptsize\CornellClauseRef{6.7.5}} & Tracks storage and object usability for Storage duration; storage availability alone does not prove that an object's lifetime has begun or remains active.\\\hline
\textbf{What rule applies to Alignment?}\\[-1mm]{\scriptsize\CornellClauseRef{6.7.6}} & Defines the core-language rules for Alignment, including formation, semantic effect, interactions with type and lookup, and any ill-formed or implementation-dependent cases.\\\hline
\textbf{What rule applies to Temporary objects?}\\[-1mm]{\scriptsize\CornellClauseRef{6.7.7}} & Defines the core-language rules for Temporary objects, including formation, semantic effect, interactions with type and lookup, and any ill-formed or implementation-dependent cases.\\\hline
\end{longtable}
\begin{CornellReferenceSummaryBox}Memory and objects supplies the core-language semantics portion of this clause. The key review move is to connect its local wording to the clause-wide model, then classify each condition by normative force before relying on it.\end{CornellReferenceSummaryBox}
\Needspace{5\baselineskip}\paragraph{Detailed coverage ledger.}
\begin{description}[style=nextline,leftmargin=2.8cm,font=\normalfont\bfseries\color{CornellReferenceTeal}]
\item[6.7.5.1] \textbf{General.} Frames the terminology, applicability, and relationships needed to interpret General; detailed rules remain in the following subclauses.
\item[6.7.5.2] \textbf{Static storage duration.} Tracks storage and object usability for Static storage duration; storage availability alone does not prove that an object's lifetime has begun or remains active.
\item[6.7.5.3] \textbf{Thread storage duration.} Tracks storage and object usability for Thread storage duration; storage availability alone does not prove that an object's lifetime has begun or remains active.
\item[6.7.5.4] \textbf{Automatic storage duration.} Tracks storage and object usability for Automatic storage duration; storage availability alone does not prove that an object's lifetime has begun or remains active.
\item[6.7.5.5] \textbf{Dynamic storage duration.} Tracks storage and object usability for Dynamic storage duration; storage availability alone does not prove that an object's lifetime has begun or remains active.
\item[6.7.5.5.1] \textbf{General.} Frames the terminology, applicability, and relationships needed to interpret General; detailed rules remain in the following subclauses.
\item[6.7.5.5.2] \textbf{Allocation functions.} Defines the core-language rules for Allocation functions, including formation, semantic effect, interactions with type and lookup, and any ill-formed or implementation-dependent cases.
\item[6.7.5.5.3] \textbf{Deallocation functions.} Defines the core-language rules for Deallocation functions, including formation, semantic effect, interactions with type and lookup, and any ill-formed or implementation-dependent cases.
\item[6.7.5.6] \textbf{Duration of subobjects.} Places Duration of subobjects in the time model; keep tick period, epoch, clock domain, civil representation, zone conversion, and ambiguous local time distinct.
\end{description}
\Needspace{8\baselineskip}
\subsection*{6.8 Types}
\begin{longtable}{@{}>{\RaggedRight\arraybackslash\columncolor{CornellReferenceCueGray}}p{0.28\textwidth}>{\RaggedRight\arraybackslash}p{0.64\textwidth}@{}}
\rowcolor{CornellReferenceNavy}\color{white}\textbf{Cue / recall prompt} & \color{white}\textbf{Notes}\\\endfirsthead
\rowcolor{CornellReferenceNavy}\color{white}\textbf{Cue / recall prompt} & \color{white}\textbf{Notes (continued)}\\\endhead
\arrayrulecolor{CornellReferenceLineGray}
\textbf{What rule applies to General?}\\[-1mm]{\scriptsize\CornellClauseRef{6.8.1}} & Frames the terminology, applicability, and relationships needed to interpret General; detailed rules remain in the following subclauses.\\\hline
\textbf{What rule applies to Fundamental types?}\\[-1mm]{\scriptsize\CornellClauseRef{6.8.2}} & Defines the core-language rules for Fundamental types, including formation, semantic effect, interactions with type and lookup, and any ill-formed or implementation-dependent cases.\\\hline
\textbf{What rule applies to Optional extended floating-point types?}\\[-1mm]{\scriptsize\CornellClauseRef{6.8.3}} & Defines the core-language rules for Optional extended floating-point types, including formation, semantic effect, interactions with type and lookup, and any ill-formed or implementation-dependent cases.\\\hline
\textbf{What rule applies to Compound types?}\\[-1mm]{\scriptsize\CornellClauseRef{6.8.4}} & Defines the core-language rules for Compound types, including formation, semantic effect, interactions with type and lookup, and any ill-formed or implementation-dependent cases.\\\hline
\textbf{What rule applies to CV-qualifiers?}\\[-1mm]{\scriptsize\CornellClauseRef{6.8.5}} & Defines the core-language rules for CV-qualifiers, including formation, semantic effect, interactions with type and lookup, and any ill-formed or implementation-dependent cases.\\\hline
\textbf{When is Conversion ranks permitted?}\\[-1mm]{\scriptsize\CornellClauseRef{6.8.6}} & Specifies source and destination conditions for Conversion ranks, the resulting type or value category, and any information-loss, pointer, qualification, or lifetime consequence.\\\hline
\end{longtable}
\begin{CornellReferenceSummaryBox}Types supplies the core-language semantics portion of this clause. The key review move is to connect its local wording to the clause-wide model, then classify each condition by normative force before relying on it.\end{CornellReferenceSummaryBox}
\Needspace{8\baselineskip}
\subsection*{6.9 Program execution}
\begin{longtable}{@{}>{\RaggedRight\arraybackslash\columncolor{CornellReferenceCueGray}}p{0.28\textwidth}>{\RaggedRight\arraybackslash}p{0.64\textwidth}@{}}
\rowcolor{CornellReferenceNavy}\color{white}\textbf{Cue / recall prompt} & \color{white}\textbf{Notes}\\\endfirsthead
\rowcolor{CornellReferenceNavy}\color{white}\textbf{Cue / recall prompt} & \color{white}\textbf{Notes (continued)}\\\endhead
\arrayrulecolor{CornellReferenceLineGray}
\textbf{What rule applies to Sequential execution?}\\[-1mm]{\scriptsize\CornellClauseRef{6.9.1}} & Defines the core-language rules for Sequential execution, including formation, semantic effect, interactions with type and lookup, and any ill-formed or implementation-dependent cases.\\\hline
\textbf{What rule applies to Multi-threaded executions and data races?}\\[-1mm]{\scriptsize\CornellClauseRef{6.9.2}} & Defines the core-language rules for Multi-threaded executions and data races, including formation, semantic effect, interactions with type and lookup, and any ill-formed or implementation-dependent cases.\\\hline
\textbf{What rule applies to Start and termination?}\\[-1mm]{\scriptsize\CornellClauseRef{6.9.3}} & Defines the core-language rules for Start and termination, including formation, semantic effect, interactions with type and lookup, and any ill-formed or implementation-dependent cases.\\\hline
\end{longtable}
\begin{CornellReferenceSummaryBox}Program execution supplies the core-language semantics portion of this clause. The key review move is to connect its local wording to the clause-wide model, then classify each condition by normative force before relying on it.\end{CornellReferenceSummaryBox}
\Needspace{5\baselineskip}\paragraph{Detailed coverage ledger.}
\begin{description}[style=nextline,leftmargin=2.8cm,font=\normalfont\bfseries\color{CornellReferenceTeal}]
\item[6.9.2.1] \textbf{General.} Frames the terminology, applicability, and relationships needed to interpret General; detailed rules remain in the following subclauses.
\item[6.9.2.2] \textbf{Data races.} Defines the core-language rules for Data races, including formation, semantic effect, interactions with type and lookup, and any ill-formed or implementation-dependent cases.
\item[6.9.2.3] \textbf{Forward progress.} Defines the core-language rules for Forward progress, including formation, semantic effect, interactions with type and lookup, and any ill-formed or implementation-dependent cases.
\item[6.9.3.1] \textbf{main function.} Defines the core-language rules for main function, including formation, semantic effect, interactions with type and lookup, and any ill-formed or implementation-dependent cases.
\item[6.9.3.2] \textbf{Static initialization.} Determines the initialization form used by Static initialization, the candidates or conversions considered, object-lifetime start, narrowing rules, and failure consequences.
\item[6.9.3.3] \textbf{Dynamic initialization of non-block variables.} Defines the blocking and synchronization contract for Dynamic initialization of non-block variables; ownership, wake conditions, timeout behavior, and happens-before edges must be analyzed together.
\item[6.9.3.4] \textbf{Termination.} Defines the core-language rules for Termination, including formation, semantic effect, interactions with type and lookup, and any ill-formed or implementation-dependent cases.
\end{description}
\section{Language-Lawyer and Library-Usage Pitfalls}
\begin{CornellReferencePitfallBox}\begin{itemize}[label=\textcolor{CornellReferenceAmber}{$\blacktriangleright$}]
\item Confusing scope, visibility, linkage, and reachability.
\item Using storage after an object's lifetime has ended.
\item Reading an indeterminate value without checking the narrow permitted cases.
\item Assuming evaluation order that the rules do not establish.
\item Treating atomicity as equivalent to ordering or synchronization.
\item Treating a note, example, or recommended practice as though it imposed a mandatory program rule.
\end{itemize}\end{CornellReferencePitfallBox}
\section{Programmer and Implementation Checklist}
\begin{CornellReferenceChecklistBox}\begin{itemize}[label=$\square$]
\item Find every declaration and determine whether it is also a definition.
\item Apply the one-definition rule across all translation units.
\item Resolve names using the correct lookup form and point of declaration.
\item Track storage duration separately from object lifetime.
\item Verify type, alignment, and object-representation assumptions.
\item Construct the happens-before argument for shared accesses.
\item Audit initialization and termination ordering.
\item For \CornellClauseRef{6.1}, verify preamble against the precise rule category and all cited dependencies.
\item For \CornellClauseRef{6.2}, verify declarations and definitions against the precise rule category and all cited dependencies.
\item For \CornellClauseRef{6.3}, verify one-definition rule against the precise rule category and all cited dependencies.
\item For \CornellClauseRef{6.4}, verify scope against the precise rule category and all cited dependencies.
\end{itemize}\end{CornellReferenceChecklistBox}
\section{Clause Synthesis}
\begin{CornellReferenceOrientationBox}Builds the core semantic model for declarations, definitions, scope, lookup, linkage, objects, lifetimes, types, execution, initialization, threading, and program start and termination. Acts as the semantic backbone of the core language. A sound review distinguishes syntax and participation from runtime preconditions, keeps notes and examples non-mandatory, and follows cross-clause dependencies before making a portability claim.\end{CornellReferenceOrientationBox}
\section{Self-Test Questions}
\begin{enumerate}
\item What role does Preamble play in the clause's overall model?
\item Which normative distinctions must be kept separate when applying Declarations and definitions?
\item How would you audit a program or implementation that depends on One-definition rule?
\item Compare Scope with Name lookup; what different question does each answer?
\item What role does Name lookup play in the clause's overall model?
\item Which normative distinctions must be kept separate when applying Program and linkage?
\item How would you audit a program or implementation that depends on Memory and objects?
\item Compare Types with Program execution; what different question does each answer?
\item What role does Program execution play in the clause's overall model?
\item Which normative distinctions must be kept separate when applying General?
\item How would you audit a program or implementation that depends on Point of declaration?
\item Compare Block scope with Function parameter scope; what different question does each answer?
\item What role does Function parameter scope play in the clause's overall model?
\item Which normative distinctions must be kept separate when applying Lambda scope?
\item Scenario: a program relies on an unstated implementation behavior while using Namespace scope. What must be checked before calling the program portable?
\item Scenario: a program relies on an unstated implementation behavior while using Class scope. What must be checked before calling the program portable?
\end{enumerate}
\clearpage\section{Answer Key}
\begin{enumerate}
\item Preamble is one part of the clause's core-language semantics model. Its role is understood by connecting the local rule in 6.1 to the clause orientation and its dependent concepts.
\item Keep formation and well-formedness separate from runtime preconditions and effects. Also distinguish mandatory wording from notes, implementation choice from undefined behavior, and interface declarations from detailed contracts; see 6.2.
\item Audit One-definition rule by locating the controlling wording in 6.3, listing inputs and type constraints, proving any preconditions, recording effects and failure behavior, and checking lifetime, invalidation, complexity, and synchronization where applicable.
\item Scope and Name lookup occupy different steps or facility groups. Analyze each under its own subclause, then follow the explicit dependency between them rather than transferring one set of guarantees to the other; begin at 6.4.
\item Name lookup is one part of the clause's core-language semantics model. Its role is understood by connecting the local rule in 6.5 to the clause orientation and its dependent concepts.
\item Keep formation and well-formedness separate from runtime preconditions and effects. Also distinguish mandatory wording from notes, implementation choice from undefined behavior, and interface declarations from detailed contracts; see 6.6.
\item Audit Memory and objects by locating the controlling wording in 6.7, listing inputs and type constraints, proving any preconditions, recording effects and failure behavior, and checking lifetime, invalidation, complexity, and synchronization where applicable.
\item Types and Program execution occupy different steps or facility groups. Analyze each under its own subclause, then follow the explicit dependency between them rather than transferring one set of guarantees to the other; begin at 6.8.
\item Program execution is one part of the clause's core-language semantics model. Its role is understood by connecting the local rule in 6.9 to the clause orientation and its dependent concepts.
\item Keep formation and well-formedness separate from runtime preconditions and effects. Also distinguish mandatory wording from notes, implementation choice from undefined behavior, and interface declarations from detailed contracts; see 6.9.2.1.
\item Audit Point of declaration by locating the controlling wording in 6.4.2, listing inputs and type constraints, proving any preconditions, recording effects and failure behavior, and checking lifetime, invalidation, complexity, and synchronization where applicable.
\item Block scope and Function parameter scope occupy different steps or facility groups. Analyze each under its own subclause, then follow the explicit dependency between them rather than transferring one set of guarantees to the other; begin at 6.4.3.
\item Function parameter scope is one part of the clause's operation semantics model. Its role is understood by connecting the local rule in 6.4.4 to the clause orientation and its dependent concepts.
\item Keep formation and well-formedness separate from runtime preconditions and effects. Also distinguish mandatory wording from notes, implementation choice from undefined behavior, and interface declarations from detailed contracts; see 6.4.5.
\item First identify the exact requirement in 6.4.6, then classify the relied-on behavior as required, unspecified, implementation-defined, conditionally supported, or undefined. Portability requires satisfying the program obligations and avoiding dependence on a choice not guaranteed in the target environments.
\item First identify the exact requirement in 6.4.7, then classify the relied-on behavior as required, unspecified, implementation-defined, conditionally supported, or undefined. Portability requires satisfying the program obligations and avoiding dependence on a choice not guaranteed in the target environments.
\end{enumerate}
\section{Key-Term Recap}
\begin{longtable}{@{}>{\RaggedRight\arraybackslash}p{0.28\textwidth}>{\RaggedRight\arraybackslash}p{0.64\textwidth}@{}}\toprule\textbf{Term} & \textbf{Working meaning}\\\midrule\endfirsthead\toprule\textbf{Term} & \textbf{Working meaning (continued)}\\\midrule\endhead\bottomrule\endfoot
definition & A declaration that supplies the entity-defining information required by the applicable rule.\\
one-definition rule & The constraints governing how many definitions may exist and agree across a program.\\
scope & The region in which a declaration can make a name usable.\\
linkage & The rule connecting declarations to the same entity across scopes or translation units.\\
object lifetime & The interval in which storage contains an object and its properties may be used under the applicable rules.\\
storage duration & The category controlling the minimum persistence of storage.\\
data race & A conflicting unsynchronized execution condition that produces undefined behavior.\\
forward progress & The guarantees governing whether threads eventually execute or complete relevant steps.\\
\end{longtable}
\CornellTakeaway{Correct C++ reasoning depends on tracking identity, lookup, type, storage, lifetime, and execution order as separate but interacting systems.}
\end{document}
